% !TeX program = xelatex
\documentclass[11pt,a4paper]{ctexart}
% ============================================================
%  网络资料对标加练 · 30 题（2026-11-08 晚）
%  指令：用户「请去网上查找一些资料或者学习视频，提供更多试题」
%  题源：kmath 交大期中卷 / MIT OCW 18.01 / LibreTexts 2.5 /
%        APEX 1.2 / Paul's 2.10 —— 只借「卡点结构」，全部换新题面与数值
%  工作流：先写 Homework/_review/tools/verify_web_drill.py（PASS 78 / FAIL 0）
%          锁定答案，再写本讲义；成文前脚本抓出 W2.3 对数求导自写错误 1 处
%  定位：11/9 早段顺延弹药，不塞进 11/8 今晚 120 分钟；不改 61 页阅读包
% ============================================================
\usepackage[a4paper,top=1.9cm,bottom=1.9cm,left=1.9cm,right=1.9cm]{geometry}
\usepackage{amsmath,amssymb,amsthm,mathtools,bm}
\usepackage{booktabs,array,longtable,multicol}
\usepackage{enumitem}
\usepackage{xcolor}
\usepackage{fancyhdr}
\usepackage{hyperref}
\hypersetup{hidelinks,colorlinks=true,linkcolor=blue!50!black,urlcolor=blue!45!black}
\usepackage{xurl}
\usepackage{titlesec}

\linespread{1.18}
\setlength{\parskip}{0.14em}
\setlist[enumerate]{itemsep=2pt,topsep=3pt,parsep=0pt,leftmargin=1.9em}
\setlist[itemize]{itemsep=2pt,topsep=3pt,parsep=0pt,leftmargin=1.6em}
\titleformat{\section}{\Large\bfseries\color{blue!35!black}}{\thesection}{0.6em}{}
\titleformat{\subsection}{\large\bfseries\color{blue!25!black}}{\thesubsection}{0.5em}{}
\titlespacing*{\section}{0pt}{1.3ex plus .3ex}{0.7ex}
\titlespacing*{\subsection}{0pt}{0.9ex plus .2ex}{0.4ex}

\pagestyle{fancy}
\fancyhf{}
\fancyhead[L]{\small 网络对标加练 · 30 题 · 11/8}
\fancyhead[R]{\small 陶理}
\fancyfoot[C]{\small\thepage}
\renewcommand{\headrulewidth}{0.4pt}

\newcommand{\R}{\mathbb{R}}
\newcommand{\eps}{\varepsilon}
\newcommand{\ee}{\mathrm{e}}
\newcommand{\D}{\mathrm{d}}
\newcommand{\Q}[2]{\par\medskip\noindent\textbf{#1}\quad #2\par}
\newcommand{\key}[1]{\par\nopagebreak\noindent{\small\color{red!55!black}\textbf{卡点：}#1}\par}
\newcommand{\tagline}[1]{\par\nopagebreak\noindent{\small\color{blue!45!black}#1}\par}

\begin{document}

\begin{center}
{\LARGE\bfseries 网络资料对标加练}\\[0.35em]
{\large 30 题 · 2026-11-08 晚}\\[0.45em]
{\small W1 极限计算 6 $+$ W2 泰勒与高阶导 6 $+$ W3 定义证明 6
 $+$ W4 导数与微分 6 $+$ W5 作图与压轴 6}
\end{center}

\begin{center}
\fcolorbox{blue!40!black}{blue!4}{
\begin{minipage}{0.95\textwidth}
\vspace{0.25em}
\textbf{怎么用}\quad 本册是 \textbf{11/9 早段顺延弹药}，\textbf{不追加到 11/8 今晚 120 分钟}（今晚按《今晚执行卡\_11-08》走）。按「题量优先」分 5 组全刷，组末立刻对 §6 速查表，正确率 $<80\%$ 回炉 3 道同型。每题标「对标」= 卡点来源，题面与数值全部新造，与本库 151 题逐题比对无重复。
\vspace{0.25em}
\end{minipage}}
\end{center}

\section{网上资料与学习视频清单（2026-11-08 检索）}

{\footnotesize
\begin{longtable}{@{}p{2.7cm}p{4.9cm}p{6.0cm}p{1.5cm}@{}}
\toprule
\textbf{类型} & \textbf{资料与地址} & \textbf{对着哪一块用} & \textbf{核实} \\
\midrule
\endfirsthead
\toprule
\textbf{类型} & \textbf{资料与地址} & \textbf{对着哪一块用} & \textbf{核实} \\
\midrule
\endhead
\bottomrule
\endfoot

交大卷（题源） & \url{https://kmath.cn/math/papername.aspx?paperid=2584} & 18 题全文：Q12 幂指首项抵消、Q14 截断多项式高阶导、Q16 作图（附 $y',y''$）、Q17 $\eps$-X 定义、Q18 中值定理两问压轴 & 已取全文 \\[2pt]

交大卷（题源） & \url{https://kmath.cn/math/papername.aspx?paperid=2430} & Q9 $\ee^x\sin x$ 高阶导、Q12 隐函数二阶导、Q13 反函数二阶导、Q14 分段最值、Q4 迭代单调 & 搜索摘要 \\[2pt]

交大卷（存档） & \url{https://www.aizhugong.com/thread-31586-1-1.html} & 2017--2019 级 A 类期中/期末 PDF 清单；需注册与金币，本次未下载 & 列表可见 \\[2pt]

题库 + 答案 & \url{https://ocw.mit.edu/courses/18-01sc-single-variable-calculus-fall-2010/download/} & MIT 18.01SC：exam1--4 与全部 solutions、pset1--8 solutions、problem-solving videos，适合「整卷手感」加练 & 已核实页面 \\[2pt]

讲义 PDF & \url{https://ocw.mit.edu/courses/18-01-single-variable-calculus-fall-2006/resources/lecture-notes/} & lec4 高阶导、lec9 线性/二次近似、lec10 Curve Sketching、lec14 MVT and Inequalities、lec37 Taylor Series & 已核实直链 \\[2pt]

视频 & MIT 18.01SC Session 4「Limits and Continuity」：Lecture 2 的 Limits / Continuity 两段 clip + Recitation「Smoothing a Piecewise Function」 & 分段函数光滑拼接、连续但不可导；英文，单段 4--9 min & 搜索摘要 \\[2pt]

定义讲义 + 习题 & \url{https://math.libretexts.org/Bookshelves/Calculus} 内 OpenStax 2.5「The Precise Definition of a Limit」及 2.5E Exercises & 单侧极限、无穷大定义、「用精确定义证明极限不存在」；2.5E 附答案 & 已核实正文 \\[2pt]

定义习题 & \url{https://spot.pcc.edu/math/APEX/sec_limit_def.html}（APEX Calculus §1.2） & 10 道 prove-the-limit（从线性到 $\sin x$、$1/x$），练量词行与 $\delta$ 反解 & 已核实正文 \\[2pt]

定义习题 & \url{http://tutorial-math.wip.lamar.edu/Problems/CalcI/DefnOfLimit.aspx}（Paul's §2.10） & 7 题带解答，含 $\lim_{x\to1}1/(x-1)^2=\infty$、$\lim_{x\to\infty}1/x^2=0$；入口 \url{https://tutorial.math.lamar.edu/} & 已核实正文 \\[2pt]

中文视频 & \url{https://www.bilibili.com/video/BV1CAxaeHEeH}（宋浩《高等数学》全程 2.0 版） & 导数定义 $\to$ 中值定理 $\to$ 泰勒 $\to$ 渐近线全链条，适合睡前只听不动笔 & 链接来自检索 \\[2pt]

中文图文 & \url{https://www.bilibili.com/read/cv40450828/}《微分中值定理、洛必达、泰勒、导数应用》 & 辅助函数构造思路 + 三条渐近线检查清单，配 W5 组 & 已核实摘要 \\[2pt]

可视化 & 3Blue1Brown《微积分的本质》第 10 集「泰勒级数」、可汗学院 Taylor series 全系列 & 泰勒多项式逐阶逼近的直觉；\textbf{只作概念补充，不替代动笔} & 线索，未核实具体地址 \\[2pt]

\end{longtable}}

\tagline{注意：kmath 等题库网页的数学公式层有 OCR 错乱（分母、指数常丢）。本册只借题型结构，所有函数与数值以本库 \texttt{.tex} 源 / sympy 复核为准，勿照抄网页公式。上表「核实」列区分了已取全文、仅见搜索摘要与未核实线索。}

\section{W1 极限计算与等价无穷小（6 题，建议 25 min）}

\Q{W1.1〔对标交大卷 Q12 幂指首项抵消〕}{求 $\displaystyle\lim_{x\to0}\frac{(1+x)^{3/x}-\ee^3(1-2x)}{x}$。}
\Q{W1.2〔对标 MIT lec9 二阶近似〕}{求 $\displaystyle\lim_{x\to0}\frac{\cos x-\ee^{-x^2/2}}{x^4}$。}
\Q{W1.3〔对标交大卷 Q11 抽象函数〕}{设 $f(1)=0$、$f'(1)=2$，求 $\displaystyle\lim_{x\to0}\frac{f(\cos x+x^2)}{x^2}$。}
\Q{W1.4〔对标 2023-24 真题复合展开〕}{求 $\displaystyle\lim_{x\to0}\frac{\sqrt{1+x^2}-\sqrt{1+\sin^2x}}{x^4}$。}
\Q{W1.5〔对标交大卷 Q3 数列极限〕}{求 $\displaystyle\lim_{n\to\infty}\Big(\frac{n!}{n^n}\Big)^{1/n}$（提示：Stirling）。}
\Q{W1.6〔对标 2015-16 Q12 通分后高阶〕}{求 $\displaystyle\lim_{x\to0}\Big(\frac1{x^2}-\frac1{\sin^2x}\Big)$。}

\section{W2 泰勒、截断多项式与高阶导（6 题，建议 35 min）}

\Q{W2.1〔对标 2023-24 Q15 $(1+x)^\alpha$ 通项〕}{设 $f(x)=(1+x)^{1/3}$，求 $f^{(9)}(0)$（写成连乘积与化简分数）。}
\Q{W2.2〔对标交大卷 Q9 $\ee^x\sin x$ 七阶导〕}{设 $f(x)=\ee^x\sin x$，求 $f^{(9)}(0)$，并给出 $f^{(n)}(0)$ 的通项。}
\Q{W2.3〔对标交大卷 Q10 对数求导〕}{设 $y=(1+x^2)^{\arctan x}$，求 $y'(1)$。}
\Q{W2.4〔对标交大卷 Q14 压轴·截断多项式〕}{记 $P_n(x)=\sum_{k=0}^{n}(-1)^k\dfrac{x^{2k}}{(2k)!}$（$n$ 为非负整数），求
$\displaystyle\left.\frac{\D^{2n+1}}{\D x^{2n+1}}\big[P_n(x)\sin x\big]\right|_{x=0}$。}
\Q{W2.5〔对标册三 P2.9 隐函数级数〕}{设 $y=y(x)$ 由 $y=x\ee^{y}$、$y(0)=0$ 确定，用形式幂级数求 $y'''(0)$。}
\Q{W2.6〔对标 MIT lec14 余项口径〕}{求 $\displaystyle\lim_{x\to0}\frac{\ln(1+x)-x+x^2/2}{x^3}$；并证明 $x\in[0,1]$ 时
$\ee^x-1-x-\frac{x^2}{2}\le\frac{\ee x^3}{6}$。说明 $x<0$ 时该不等式是否仍成立。}

\section{W3 极限定义证明（6 题，建议 35 min）}

\Q{W3.1〔$\eps$-$\delta$；对标 APEX §1.2 习题〕}{用定义证明 $\lim_{x\to2}x^3=8$，写出可用的 $\delta(\eps)$ 并验证。}
\Q{W3.2〔$\eps$-$\delta$ + 有理化；对标 2011-12 Q11〕}{用定义证明 $\lim_{x\to0}\dfrac{\sqrt{1+x}-1}{x}=\dfrac12$，写出可用的 $\delta(\eps)$。}
\Q{W3.3〔$\eps$-X + 有界振荡；对标 2019-20 Q11〕}{用定义证明 $\lim_{x\to+\infty}\dfrac{3x+\sin x}{x+1}=3$。}
\Q{W3.4〔$\eps$-X + 对数放缩；对标 MIT 18.01SC〕}{用定义证明 $\lim_{x\to+\infty}\dfrac{\ln(1+x)}{x}=0$（提示：$\ln(1+x)\le2\sqrt{x}$，$x\ge1$）。}
\Q{W3.5〔M 型；对标 2024-25 Q11 放缩陷阱〕}{用定义证明 $\lim_{x\to0}\dfrac1{|x|(2+|x|)}=\infty$，写出 $\delta(M)$。注意只能写「$|f|\to\infty$」。}
\Q{W3.6〔定义证「不存在」；对标 LibreTexts 2.5E〕}{用精确定义证明 $\lim_{x\to0}\sin\frac1x$ 不存在（给出两列趋于 0 的点）。}

\section{W4 导数、隐函数、反函数与极坐标（6 题，建议 25 min）}

\Q{W4.1〔对标交大卷 Q6 函数方程〕}{设 $3f(2x)-f(1-2x)=4x-2$ 对一切 $x$ 成立，求 $f(x)$。}
\Q{W4.2〔对标交大卷 Q13 反函数二阶导〕}{设 $y=x+\ln(1+x)$，其反函数为 $x=f(y)$，求 $f'(0)$ 与 $f''(0)$。}
\Q{W4.3〔对标交大卷 Q12 隐函数二阶导〕}{设 $y=y(x)$ 由 $x\ee^{y}+y\ee^{x}=0$、$y(0)=0$ 确定，求 $y'(0)$ 与 $y''(0)$。}
\Q{W4.4〔对标 2024-25 Q8 极坐标切线〕}{求曲线 $r=\sin2\theta$ 在 $\theta=\pi/6$ 处的切线斜率（$0/0$ 情形须取极限）。}
\Q{W4.5〔对标交大卷 Q13 分段二阶可导〕}{设 $f(x)=\begin{cases}\dfrac{\sin x}{x},&x>0\\[2pt] ax^2+bx+c,&x\le0\end{cases}$ 二阶可导，求 $a,b,c$ 及 $f''(0)$。}
\Q{W4.6〔对标交大卷 Q8 参数方程〕}{设 $x=t^2-1,\ y=t^3+t$，求 $\left.\dfrac{\D^2y}{\D x^2}\right|_{t=1}$。}

\newpage
\section{W5 作图、Rolle 构造、函数方程与迭代（6 题，建议 45 min）}

\Q{W5.1〔对标交大卷 Q16 作图大题，10--12 分〕}{按七步全面讨论 $y=\dfrac{x^2}{x-1}$ 的性态并作图：定义域、奇偶与周期、$y'$ 与 $y''$、单调与极值、凹凸与拐点、\textbf{三类渐近线两端都查}、关键点列表。}
\Q{W5.2〔对标交大卷 Q18(1) 中值点〕}{设 $f$ 在 $[0,1]$ 连续、$(0,1)$ 可导，且 $f(0)=f(1)=0$。证明存在 $\eta\in(0,1)$ 使 $\eta f'(\eta)+2f(\eta)=0$。}
\Q{W5.3〔对标交大卷 Q18(2) 二阶压轴〕}{设 $f$ 在 $[0,1]$ 上二阶可导且 $f(0)=f(1)=0$。证明存在 $\xi\in(0,1)$ 使
$\xi^2f''(\xi)+4\xi f'(\xi)+2f(\xi)=0$。}
\Q{W5.4〔对标 2024-25 Q18 抽象函数方程〕}{设 $f$ 在 $\R$ 上连续，且 $f(x+y)=f(x)+f(y)+xy(x+y)$、$f(1)=1$。求 $f(x)$ 并说明解唯一。}
\Q{W5.5〔对标交大卷 Q4 + 册三 P3 迭代〕}{设 $a_1=1,\ a_{n+1}=\ln(1+a_n)$。证明 $a_n>0$ 且严格递减、$\lim a_n=0$，并求 $\lim_{n\to\infty}n a_n$。}
\Q{W5.6〔对标 MIT lec14 MVT and Inequalities〕}{证明 $x\in[0,1]$ 时 $\ee^x\ge1+x+\dfrac{x^2}{2}$；再判断 $x<0$ 时是否仍成立，给出理由或反例。}

\section{答案速查（组末立刻对，只记题号与错因）}

{\footnotesize
\begin{longtable}{@{}|c|p{5.0cm}|p{9.0cm}|@{}}
\hline
题号 & 答案 & 一句话卡点 \\
\hline
\endfirsthead
\hline
题号 & 答案 & 一句话卡点 \\
\hline
\endhead
\hline
\endfoot
W1.1 & $\ee^3/2$ & $(1+x)^{3/x}=\ee^3\ee^{-3x/2+O(x^2)}$，首项抵消后看次项 \\
W1.2 & $-1/12$ & 两者二阶项相同，必须展到 $x^4$，等价无穷小失效 \\
W1.3 & $1$ & $f(1+t)=2t+o(t)$，$\cos x+x^2-1=x^2/2+O(x^4)$ \\
W1.4 & $1/6$ & $\sin^2x=x^2-x^4/3+O(x^6)$ 后再开方 \\
W1.5 & $1/\ee$ & Stirling 取对数，或比值法 $\frac{a_{n+1}}{a_n}=(1+\frac1n)^{-n}\to\ee^{-1}$ \\
W1.6 & $-1/3$ & 通分后分子需展到 $x^6$，别在 $x^4$ 就停 \\
W2.1 & $\prod_{k=0}^{8}(\frac13-k)=\frac{96342400}{19683}$ & $[x^9]$ 系数乘 $9!$；通项导数\textbf{不}再乘 $9!$ \\
W2.2 & $16$；$f^{(n)}(0)=2^{n/2}\sin\frac{n\pi}{4}$ & 取 $\operatorname{Im}\ee^{(1+i)x}$，$n=9$ 给 $16$ \\
W2.3 & $2^{\pi/4}\big(\frac\pi4+\frac{\ln2}{2}\big)$ & 对数求导两项：$2x\arctan x/(1+x^2)$ 与 $\ln(1+x^2)/(1+x^2)$ \\
W2.4 & $(-4)^n$ & $[x^{2n+1}]\Rightarrow(-1)^n\sum_{k=0}^{n}\binom{2n+1}{2k}=(-1)^n4^n$ \\
W2.5 & $9$ & $y=x+x^2+\frac32x^3+\cdots$，故 $y'''(0)=6\cdot\frac32$ \\
W2.6 & $1/3$；Lagrange 余项成立 & Peano 只给阶数；Lagrange 才给符号与上界；$x=-1$ 反向 \\
W3.1 & $\delta=\min\{1,\eps/19\}$ & 先限 $|x-2|<1$，再界住 $x^2+2x+4\le19$ \\
W3.2 & $\delta=\min\{1/2,2\eps\}$ & 有理化 $f=1/(\sqrt{1+x}+1)$，差 $\le|x|/5$ \\
W3.3 & $X=\max\{1,4/\eps\}$ & 差 $=\frac{\sin x-3}{x+1}$，分子 $\le4$ \\
W3.4 & $X=\max\{1,4/\eps^2\}$ & $\ln(1+x)\le2\sqrt x\Rightarrow$ 差 $\le2/\sqrt x$ \\
W3.5 & $\delta=\min\{\frac12,\frac{2}{5M}\}$ & 先限 $|x|\le1/2$ 得 $|x|(2+|x|)\le\frac52|x|$，\textbf{分母要上界} \\
W3.6 & 不存在 & $x_k=\frac1{2k\pi+\pi/2}$ 与 $y_k=\frac1{2k\pi+3\pi/2}$，值恒为 $\pm1$ \\
W4.1 & $f(x)=\frac x2-\frac34$ & $u=2x$ 后再以 $1-u$ 代入，解二元方程组 \\
W4.2 & $f'(0)=\frac12,\ f''(0)=\frac18$ & $f''=-g''/(g')^3$，$g'(0)=2,\ g''(0)=-1$ \\
W4.3 & $y'(0)=-1,\ y''(0)=4$ & $F_{xy}(0,0)=2$，二阶公式 $y''=-\frac{F_{xx}+2F_{xy}y'+F_{yy}y'^2}{F_y}$ \\
W4.4 & $\frac{5\sqrt3}{3}$ & $\frac{r'\sin\theta+r\cos\theta}{r'\cos\theta-r\sin\theta}$，$\theta=\pi/6$ 时 $r=\frac{\sqrt3}{2},r'=1$ \\
W4.5 & $a=-\frac16,\ b=0,\ c=1$；$f''(0)=-\frac13$ & $\frac{\sin x}{x}=1-\frac{x^2}{6}+O(x^4)$ 逐项对齐 \\
W4.6 & $\frac12$ & $\frac{\D}{\D t}\big(\frac{y'_t}{x'_t}\big)/x'_t$，别漏再除一次 $x'_t$ \\
W5.1 & $y'=\frac{x(x-2)}{(x-1)^2},\ y''=\frac2{(x-1)^3}$ & 垂直 $x=1$（左 $-\infty$ 右 $+\infty$）、斜 $y=x+1$；极大 $(0,0)$、极小 $(2,4)$、无拐点 \\
W5.2 & 构造 $F=x^2f$ & $F(0)=F(1)=0\Rightarrow F'(\eta)=0\Rightarrow 2f(\eta)+\eta f'(\eta)=0$ \\
W5.3 & 构造 $F=x^2f$，两次 Rolle & $F'(0)=0$ 且 $\exists c:F'(c)=0\Rightarrow\exists\xi:F''(\xi)=0$ \\
W5.4 & $f(x)=\frac{x^3}{3}+\frac{2x}{3}$ & 令 $g=f-\frac{x^3}{3}$ 化 Cauchy 可加，连续 $\Rightarrow g(x)=\frac23x$ \\
W5.5 & $\lim a_n=0$；$\lim n a_n=2$ & $\frac1{a_{n+1}}-\frac1{a_n}\to\frac12$（展开 $\ln(1+a)=a-\frac{a^2}{2}+\cdots$） \\
W5.6 & $x\in[0,1]$ 成立；$x<0$ 不成立 & $R_2=\frac{\ee^\xi}{6}x^3\ge0$；$\ee^{-1}\approx0.368<0.5$ \\
\hline
\end{longtable}}

\section{逐题详解（卡住满 5 min 才看，合上重写）}

\Q{W1.1}{$\frac{3\ln(1+x)}{x}=3-\frac{3x}{2}+O(x^2)$，故 $(1+x)^{3/x}=\ee^3\big(1-\frac{3x}{2}+O(x^2)\big)$；分子 $=\ee^3\big[(1-\frac{3x}{2})-(1-2x)\big]+O(x^2)=\frac{\ee^3}{2}x+O(x^2)$。}
\key{幂指函数不能直接代入 $\ee^3$ 就判 $0/0$；必须先取对数展到 $x^1$。}

\Q{W1.2}{$\cos x=1-\frac{x^2}{2}+\frac{x^4}{24}+O(x^6)$，$\ee^{-x^2/2}=1-\frac{x^2}{2}+\frac{x^4}{8}+O(x^6)$，差 $=\big(\frac1{24}-\frac18\big)x^4=-\frac{x^4}{12}$。}
\key{二阶项一模一样，$\cos x-1\sim-x^2/2$ 这类等价无穷小在这里\textbf{全部失效}。}

\Q{W1.3}{$f(1+t)=f(1)+f'(1)t+o(t)=2t+o(t)$；$\cos x+x^2-1=\frac{x^2}{2}+O(x^4)$，故分子 $=x^2+o(x^2)$。}
\key{抽象函数题只许用「一点可导」给的一阶信息，不要假设 $f''$ 存在。}

\Q{W1.4}{$\sin^2x=x^2-\frac{x^4}{3}+O(x^6)$；$\sqrt{1+u}=1+\frac u2-\frac{u^2}{8}+O(u^3)$。第一式 $u=x^2$，第二式 $u=x^2-\frac{x^4}{3}$，相减得 $\frac{x^4}{6}+O(x^6)$。}
\key{复合展开要保留内层到 $x^4$，否则答案会误成 $0$。}

\Q{W1.5}{令 $a_n=\frac{n!}{n^n}$。法一（Stirling）：$\frac1n\ln a_n=\frac1n\big[\ln\sqrt{2\pi n}+n\ln n-n+\frac1{12n}+o(1)-n\ln n\big]\to-1$。法二（比值法）：$\lim a_n^{1/n}=\lim\frac{a_{n+1}}{a_n}=\lim\frac{n^n}{(n+1)^n}=\lim\big(1+\frac1n\big)^{-n}=\ee^{-1}$。}
\key{比值法这里\textbf{可以用}；约分时 $(n+1)!=(n+1)n!$ 与 $(n+1)^{n+1}=(n+1)(n+1)^n$ 整体相消，别把指数上的 $n$ 与 $n+1$ 当成同一个。}

\Q{W1.6}{$\frac{\sin^2x-x^2}{x^2\sin^2x}$；$\sin x=x-\frac{x^3}{6}+\frac{x^5}{120}+O(x^7)\Rightarrow\sin^2x=x^2-\frac{x^4}{3}+\frac{2x^6}{45}+O(x^8)$，故分子 $=-\frac{x^4}{3}+\frac{2x^6}{45}$，分母 $=x^4-\frac{x^6}{3}$，商 $\to-\frac13$。}
\key{分母也要展到 $x^6$，只展分子会漏。}

\Q{W2.1}{$f^{(9)}(0)=9!\,[x^9](1+x)^{1/3}=\prod_{k=0}^{8}(\frac13-k)=\frac{96342400}{19683}$。}
\key{$\alpha(\alpha-1)\cdots(\alpha-n+1)$ 本身就是导数值；「再乘 $9!$」是本库复核脚本最早抓到的错。}

\Q{W2.2}{$\ee^x\sin x=\operatorname{Im}\ee^{(1+i)x}\Rightarrow f^{(n)}(0)=\operatorname{Im}(1+i)^n=2^{n/2}\sin\frac{n\pi}{4}$；$n=9$ 时 $2^{9/2}\cdot\frac{\sqrt2}{2}=16$。}
\key{奇函数/偶函数拆项也能做，但复数法一行到底且不易错符号。}

\Q{W2.3}{$\ln y=\arctan x\cdot\ln(1+x^2)$，$\frac{y'}{y}=\frac{\ln(1+x^2)}{1+x^2}+\frac{2x\arctan x}{1+x^2}$；$x=1$ 时 $y=2^{\pi/4}$，故 $y'(1)=2^{\pi/4}\big(\frac{\ln2}{2}+\frac\pi4\big)$。}
\key{乘积两因子都含 $x$，漏掉 $\arctan x$ 的导数项就会少 $\pi/4$。}

\Q{W2.4}{$[x^{2n+1}]\big(P_n\sin x\big)=(-1)^n\sum_{k=0}^{n}\frac{1}{(2k)!\,(2n-2k+1)!}$；乘 $(2n+1)!$ 得 $(-1)^n\sum_{k=0}^{n}\binom{2n+1}{2k}=(-1)^n4^n$，即 $(-4)^n$。偶数项二项式和 $=\frac12[(1+1)^{2n+1}+(1-1)^{2n+1}]=4^n$。}
\key{这是对标题眼：\textbf{截断}到 $2n$ 次，高于 $2n$ 的项不参与 $[x^{2n+1}]$；若把 $P_n$ 当成完整 $\cos x$ 会得到 $2^{2n}$ 之外的错值。}

\Q{W2.5}{设 $y=a_1x+a_2x^2+a_3x^3+\cdots$ 代入 $y=x\ee^y=x(1+y+\frac{y^2}{2}+\cdots)$，比较 $x,x^2,x^3$ 系数得 $a_1=1,\ a_2=1,\ a_3=\frac32$，故 $y'''(0)=6a_3=9$。}
\key{$y$ 里含 $x^3$ 时 $\ee^y$ 必须展到 $y^2$ 的 $x^3$ 项，只展一阶会得 $a_3=1$。}

\Q{W2.6}{$\ln(1+x)=x-\frac{x^2}{2}+\frac{x^3}{3}+O(x^4)$。带 Lagrange 余项 $\ee^x=1+x+\frac{x^2}{2}+\frac{\ee^\xi x^3}{6}$（$\xi$ 在 $0$ 与 $x$ 之间），$x\in[0,1]$ 时 $\ee^\xi\le\ee$，得不等式；$x=-1$ 时 $\ee^{-1}<1-1+\frac12$ 反向。}
\key{Peano 余项只说明「差是 $o(x^2)$」，\textbf{不给符号、不给上界}；涉及不等式必须换 Lagrange。}

\Q{W3.1}{$|x^3-8|=|x-2|\,|x^2+2x+4|$；先取 $\delta_1=1$ 使 $x\in(1,3)$，此时 $|x^2+2x+4|\le19$，再取 $\delta=\min\{1,\eps/19\}$。}
\key{量词行顺序：$\forall\eps>0$、$\exists\delta>0$、$0<|x-2|<\delta\Rightarrow\cdots$，三步缺一不可。}

\Q{W3.2}{有理化 $f=\frac1{\sqrt{1+x}+1}$，$\big|f-\frac12\big|=\frac{|x|}{2(\sqrt{1+x}+1)^2}$；先限 $|x|\le\frac12$ 使 $\sqrt{1+x}\ge\frac1{\sqrt2}$，分母 $\ge5$，故 $\delta=\min\{\frac12,2\eps\}$ 足够（数值已验）。}
\key{$\delta$ 不必是「最大」的，只要能闭环；但必须先用 $\delta_1$ 把 $x$ 关进小区间，否则根号下界无从谈起。}

\Q{W3.3}{$\big|\frac{3x+\sin x}{x+1}-3\big|=\frac{|\sin x-3|}{x+1}\le\frac4{x+1}$（$x>0$），取 $X=\max\{1,4/\eps\}$。}
\key{$\sin x$ 有界但不收敛，只能放成常数 $4$；不要写「$\sin x\to$ 某值」。}

\Q{W3.4}{$x\ge1$ 时 $\ln(1+x)\le2\sqrt x$（脚本已逐点验证），故 $\frac{\ln(1+x)}x\le\frac2{\sqrt x}<\eps$ 只要 $x>\frac4{\eps^2}$，取 $X=\max\{1,4/\eps^2\}$。}
\key{$\eps$-X 题里 $X$ 可以依赖 $\eps^2$；别硬套「$X=C/\eps$」的线性模板。}

\Q{W3.5}{先限 $|x|\le\frac12$，则 $|x|(2+|x|)\le\frac52|x|$，于是 $|f|\ge\frac{2}{5|x|}>M$ 只要 $|x|<\frac{2}{5M}$；$\delta=\min\{\frac12,\frac{2}{5M}\}$。}
\key{求 $|f|$ 的\textbf{下界}时分母要\textbf{上界}。2024-25 官方答案在这里写反（$\frac1{|x|(1+|x|)}\ge\frac1{|x|}$ 不成立），本库订正为 $\delta=\min\{\frac12,\frac1{2G}\}$。}

\Q{W3.6}{取 $x_k=\frac1{2k\pi+\pi/2}$、$y_k=\frac1{2k\pi+3\pi/2}$，二者都 $\to0$，但 $\sin(1/x_k)=1$、$\sin(1/y_k)=-1$。任给 $\delta>0$，取 $k$ 充分大使两点都落在 $(0,\delta)$ 内，函数值差恒为 $2$，与极限唯一性矛盾。}
\key{「不存在」也必须\textbf{用定义}证：给出两列点 + 固定 $\eps_0=\frac12$，而不是只说「振荡」。}

\Q{W4.1}{令 $u=2x$：$3f(u)-f(1-u)=2u-2$；再把 $u$ 换成 $1-u$ 得 $-f(u)+3f(1-u)=-2u$。解方程组得 $f(u)=\frac u2-\frac34$。}
\key{一次代入只有一个方程、两个未知函数值；必须「对称换元」凑成方程组。}

\Q{W4.2}{$g(x)=x+\ln(1+x)$，$g(0)=0,\ g'(0)=2,\ g''(0)=-1$；$f'(0)=\frac1{g'(0)}=\frac12$，$f''(0)=-\frac{g''(0)}{g'(0)^3}=\frac18$。}
\key{反函数二阶导公式的符号与三次方最易错。}

\Q{W4.3}{$F=x\ee^y+y\ee^x$；$F_x=\ee^y+y\ee^x$、$F_y=x\ee^y+\ee^x$，在 $(0,0)$ 均为 $1$，故 $y'=-1$。$F_{xx}(0,0)=0$、$F_{xy}(0,0)=2$、$F_{yy}(0,0)=0$，代入 $y''=-\frac{F_{xx}+2F_{xy}y'+F_{yy}y'^2}{F_y}=4$。}
\key{二阶隐函数求导必须用\textbf{全二阶公式}，不能对 $y'=-F_x/F_y$ 只做链式法则而漏 $F_{xy}$。}

\Q{W4.4}{$r=\sin2\theta,\ r'=2\cos2\theta$；$\theta=\frac\pi6$ 时 $r=\frac{\sqrt3}{2},\ r'=1$，分子 $r'\sin\theta+r\cos\theta=\frac12+\frac34=\frac54$，分母 $r'\cos\theta-r\sin\theta=\frac{\sqrt3}{2}-\frac{\sqrt3}{4}=\frac{\sqrt3}{4}$，斜率 $=\frac{5}{\sqrt3}=\frac{5\sqrt3}{3}$。}
\key{若出现 $0/0$（如 2024-25 的 $r=1+\cos3\theta$）必须对 $\theta$ 取极限，不能直接判斜率不存在。}

\Q{W4.5}{$\frac{\sin x}{x}=1-\frac{x^2}{6}+\frac{x^4}{120}+O(x^6)$；连续 $\Rightarrow c=1$，一阶导 $\Rightarrow b=0$，二阶导 $\Rightarrow 2a=-\frac13$ 即 $a=-\frac16$，$f''(0)=-\frac13$。}
\key{「二阶可导」要查到 $f''$ 连续衔接，不是只查 $f$ 与 $f'$。}

\Q{W4.6}{$\frac{\D y}{\D x}=\frac{3t^2+1}{2t}$，$\frac{\D^2y}{\D x^2}=\frac{\D}{\D t}\big(\frac{3t^2+1}{2t}\big)\cdot\frac1{2t}=\frac{3t^2-1}{4t^3}$，$t=1$ 给 $\frac12$。}
\key{参数方程二阶导 = 对 $t$ 求导\textbf{再除一次} $x'_t$。}

\Q{W5.1}{定义域 $x\ne1$；非奇非偶、非周期。$y'=\frac{x(x-2)}{(x-1)^2}$，$y''=\frac{2}{(x-1)^3}$。驻点 $x=0,2$：$(-\infty,0)$ 增、$(0,1)$ 减、$(1,2)$ 减、$(2,+\infty)$ 增；极大 $(0,0)$、极小 $(2,4)$。$x<1$ 时 $y''<0$、$x>1$ 时 $y''>0$，两侧凹凸性相反，但 $x=1$ 不在定义域内、\textbf{拐点必须是曲线上的点}，故无拐点。渐近线：$x=1$ 为垂直渐近线，$\lim_{x\to1^-}=-\infty$、$\lim_{x\to1^+}=+\infty$；$y-(x+1)=\frac1{x-1}\to0$（$x\to\pm\infty$ 都要写），故斜渐近线 $y=x+1$，无水平渐近线；无坐标轴交点以外的截距（过原点）。}
\key{七步里最容易丢的是「斜渐近线两端都查」与「垂直渐近线两侧符号」；卷面常附 $y',y''$，考的是流程不是硬算。}

\Q{W5.2}{令 $F(x)=x^2f(x)$，则 $F(0)=F(1)=0$，Rolle 给 $\eta\in(0,1)$ 使 $F'(\eta)=0$；而 $F'=\eta(2f+\eta f')$ 且 $\eta\ne0$，故 $\eta f'(\eta)+2f(\eta)=0$。}
\key{看到「$a f'+b f$」就往 $\ee^{\lambda x}x^{k}f$ 型辅助函数凑，本题是 $x^2f$。}

\Q{W5.3}{仍取 $F=x^2f$。$F(0)=F(1)=0\Rightarrow\exists c\in(0,1):F'(c)=0$；又 $F'(0)=0$（$F'=2xf+x^2f'$）。对 $F'$ 在 $[0,c]$ 用 Rolle 得 $\xi\in(0,c)$ 使 $F''(\xi)=0$。而 $F''=x^2f''+4xf'+2f$，即结论。}
\key{第二次 Rolle 的第二个零点是 $F'(0)=0$，不是 $F(0)=0$；这一句是整题的得分点。}

\Q{W5.4}{注意 $\frac{(x+y)^3-x^3-y^3}{3}=xy(x+y)$。令 $g(x)=f(x)-\frac{x^3}{3}$，则 $g(x+y)=g(x)+g(y)$；$g$ 连续的可加函数必为 $g(x)=kx$，$g(1)=1-\frac13=\frac23$，故 $f(x)=\frac{x^3}{3}+\frac{2x}{3}$。}
\key{不写「连续 $\Rightarrow$ 线性」这一步，答案只是猜出来的。}

\Q{W5.5}{$0<a_{n+1}=\ln(1+a_n)<a_n$（$a_n>0$），故 $a_n\downarrow L\ge0$；取极限 $L=\ln(1+L)$ 得 $L=0$。由 $\ln(1+a)=a-\frac{a^2}{2}+\frac{a^3}{3}-\cdots$ 得 $\frac1{a_{n+1}}-\frac1{a_n}=\frac12-\frac{a_n}{12}+O(a_n^2)\to\frac12$，Cesàro 给 $\frac1{a_n}\sim\frac n2$，故 $n a_n\to2$。}
\key{$\lim a_n=0$ 不等于 $\lim n a_n=0$；渐近速度必须用 $\frac1{a_{n+1}}-\frac1{a_n}$ 的极限。}

\Q{W5.6}{带 Lagrange 余项展开到二阶：$\ee^x=1+x+\frac{x^2}{2}+\frac{\ee^\xi x^3}{6}$（$\xi$ 在 $0,x$ 之间）。$x\ge0$ 时余项 $\ge0$；$x<0$ 时 $x^3<0$ 使余项为负，例如 $x=-1$ 时 $\ee^{-1}\approx0.3679<0.5$。}
\key{「展开式对一切 $x$ 成立」与「不等式对一切 $x$ 成立」是两件事，别把 Peano 余项当符号依据。}

\section{回炉索引与 11/9 衔接}

\begin{itemize}
\item \textbf{分组与反馈}：W1（25 min）/ W2（35 min）/ W3（35 min）/ W4（25 min）/ W5（45 min）共 165 min，分 5 组；组末对 §6，正确率 $<80\%$ 回炉 3 道同型。
\item \textbf{按卡点回炉}：展错阶 $\to$ W1.1/W1.2/W1.4/W2.6；通项导数 $\to$ W2.1/W2.2/W2.4；M 型放缩 $\to$ W3.5 + 执行卡 G4/G9；隐函数二阶 $\to$ W4.3/W4.6；辅助函数 $\to$ W5.2/W5.3；迭代渐近 $\to$ W5.5。
\item \textbf{排期}：今晚只按《今晚执行卡\_11-08》的四项走；11/9 早段先补完 11/8 顺延的 G 组，再做模拟卷 A 整卷 120 min + 60--75 min 订正，本册 30 题按剩余时间分 2--3 组插入，\textbf{订正时间不得压缩}。
\item \textbf{总量}：库内 151 题 $+$ 本册 30 题 $=181$ 题，另 11/9 A 卷 18 题、11/10 B 卷 18 题。本册\textbf{不新增整卷}、不改 61 页阅读包的绝对页码。
\item \textbf{复核}：\texttt{Homework/\_review/tools/verify\_web\_drill.py}，PASS 78 / FAIL 0；所有 $\delta(\eps)$、$X(\eps)$、$\delta(M)$ 均另做数值有效性检验。
\end{itemize}

\end{document}
